\chapter{O Sacramento do Batismo}

O Catecismo da Igreja Católica dedica ao Sacramento do Batismo os parágrafos 1213 a 1284, situando-o como o primeiro e fundamental sacramento da iniciação cristã. O Batismo é o ``fundamento de toda a vida cristã'', a porta que dá acesso aos outros sacramentos e o meio pelo qual somos libertos do pecado e regenerados como filhos de Deus. É o sacramento da fé por excelência, porque é nele que, em resposta à Palavra de Deus, cada batizado entra sacramentalmente na vida da Igreja e se incorpora ao Corpo de Cristo. Ao mesmo tempo, o Catecismo ilumina o Batismo como cumprimento de toda uma longa história de preparação, que remonta às águas primordiais da criação e percorre as grandes obras salvíficas de Deus na Antiga Aliança.

\section{O fundamento da vida cristã e seus nomes (1213--1216)}

O Sacramento do Batismo é o fundamento de toda a vida cristã, o pórtico da vida no Espírito e a porta de acesso aos outros sacramentos. Por ele somos libertos do pecado e regenerados como filhos de Deus, tornamo-nos membros de Cristo e somos incorporados na Igreja, tornando-nos participantes na sua missão. Pode ser definido, segundo a fórmula catequética tradicional, como ``o sacramento da regeneração pela água e pela Palavra''. Este ponto de partida da existência cristã tem uma riqueza tão grande que a própria Tradição lhe atribuiu vários nomes, cada um dos quais ilumina um aspecto diferente do mesmo mistério.

O nome ``Batismo'' vem do rito central com que se celebra: batizar (do grego \emph{baptizein}) significa ``mergulhar'' ou ``imergir''. A imersão na água simboliza a sepultura do catecúmeno na morte de Cristo, da qual ele emerge pela ressurreição com Ele como ``nova criatura''. Este simbolismo pascal está no coração do Batismo: o batizado não apenas recebe um sinal exterior, mas é realmente configurado com a morte e ressurreição do Senhor. Por isso, a forma mais significativa de conferir o Batismo é a tríplice imersão, que exprime plenamente este mistério de morte e vida nova.

O sacramento é também chamado ``banho de regeneração e renovação pelo Espírito Santo'', porque significa e realiza o nascimento da água e do Espírito, sem o qual ninguém pode entrar no Reino de Deus. Como ``iluminação'', o Batismo exprime o fato de que o batizado, ao receber o Verbo, ``luz verdadeira que ilumina todo o homem'', torna-se ele próprio ``filho da luz'' e luz para o mundo. Já no século IV, São Gregório Nazianzo reunia vários destes nomes numa bela síntese: ``Chamamos-lhe dom, graça, unção, iluminação, veste de incorruptibilidade, banho de regeneração, selo e tudo o que há de mais precioso''.

Cada um destes títulos ilumina um aspecto do mesmo mistério: é chamado ``dom'' porque é conferido gratuitamente; ``graça'' porque é dado até mesmo aos culpados; ``unção'' pelo seu caráter sacerdotal e real; ``veste'' porque cobre a vergonha do pecado; ``banho'' porque lava; e ``selo'' porque nos guarda como sinal do senhorio de Deus. Esta riqueza de nomes revela que o Batismo não é uma realidade unívoca e simples, mas um evento espiritual de profundidades inexauríveis, que toca a totalidade da existência humana e a orienta para Deus. Por isso, a catequese batismal nunca pode reduzir-se a uma mera explicação ritual, mas deve ser uma iniciação progressiva ao mistério da nova vida em Cristo.

\section{O Batismo na história da salvação (1217--1228)}

Na liturgia da Vigília Pascal, ao abençoar a água batismal, a Igreja faz solenemente memória dos grandes acontecimentos da história da salvação que já prefiguravam o mistério do Batismo. Esta memória não é um simples exercício arqueológico, mas uma revelação da continuidade profunda entre as obras de Deus na Antiga Aliança e o que acontece hoje em cada cristão que se aproxima da pia batismal. Desde o princípio do mundo a água, ``esta criatura humilde e admirável'', é fonte de vida e fecundidade, ``incubada'' pelo Espírito de Deus na criação. A Bênção da Água Pascal aclama: ``Logo no princípio do mundo, o vosso Espírito pairava sobre as águas, para que já desde então concebessem o poder de santificar''.

A Igreja viu na arca de Noé uma primeira prefiguração da salvação pelo Batismo: assim como ``um pequeno grupo, ao todo oito pessoas, foram salvos pela água'' do dilúvio, também o Batismo salva do dilúvio do pecado e inicia um novo começo de santidade. A passagem do Mar Vermelho é a prefiguração mais explícita: a libertação de Israel da escravidão do Egito anuncia a libertação operada pelo Batismo, que transfere o batizado do domínio do pecado para a liberdade dos filhos de Deus. Da mesma forma, a travessia do Jordão, pela qual o povo recebeu a terra prometida à descendência de Abraão, prefigura o Batismo como entrada na herança da vida eterna, dom da Nova Aliança. Assim, toda a história da salvação converge para as águas batismais.

Todas estas prefigurações encontram a sua realização em Jesus Cristo. Ele inicia a sua vida pública fazendo-se batizar por João no Jordão, submetendo-se voluntariamente a esse rito para ``cumprir toda a justiça'': o Espírito Santo desce sobre Ele como prelúdio da nova criação, e o Pai manifesta Jesus como ``Filho muito amado''. Cristo já havia falado da sua Paixão como de um ``batismo'' com que devia ser batizado; e é da sua Páscoa que brotam as fontes do Batismo cristão. O sangue e a água que manaram do lado aberto de Cristo crucificado são tipos do Batismo e da Eucaristia, sacramentos da vida nova: ``Onde é que foste batizado, de onde é que vem o Batismo, senão da cruz de Cristo, da morte de Cristo? Ali está todo o mistério''.

Desde o Pentecostes, a Igreja celebra e administra o Batismo, sempre ligado à fé e à conversão. São Pedro declara à multidão: ``Convertei-vos e peça cada um o Batismo em nome de Jesus Cristo, para vos serem perdoados os pecados; recebereis então o dom do Espírito Santo''. Segundo São Paulo, pelo Batismo o crente comunga na morte de Cristo --- é sepultado e ressuscita com Ele: ``Todos nós, que fomos batizados em Cristo Jesus, fomos batizados na sua morte... para que, assim como Cristo ressuscitou dos mortos, também nós vivamos uma vida nova''. O Batismo é, pois, um banho de água no qual ``a semente incorruptível'' da Palavra de Deus produz o seu efeito vivificador: ``Junta-se a palavra ao elemento material e faz-se o sacramento'', como disse Santo Agostinho.

\section{A celebração do Batismo (1229--1245)}

Desde os tempos apostólicos, tornar-se cristão requer um caminho e uma iniciação com diversas etapas. Este itinerário inclui elementos essenciais: o anúncio da Palavra, o acolhimento do Evangelho que implica a conversão, a profissão de fé, o Batismo, a efusão do Espírito Santo e o acesso à comunhão eucarística. Nos primeiros séculos, a iniciação cristã conheceu grande desenvolvimento, com um longo catecumenato e uma série de ritos preparatórios escalonados liturgicamente. O Concílio Vaticano II restaurou, para a Igreja latina, ``o catecumenato dos adultos, distribuído em várias fases'', cujo ritual se encontra no \emph{Ordo initiationis christianae adultorum} (1972).

Hoje em dia, em todos os ritos latinos e orientais, a iniciação cristã dos adultos começa com a entrada no catecumenato e culmina na celebração única dos três sacramentos --- Batismo, Confirmação e Eucaristia ---, preferivelmente na Vigília Pascal. O catecumenato tem por finalidade conduzir os candidatos, em resposta à iniciativa divina e em união com a comunidade eclesial, à maturidade da sua conversão e da sua fé: é uma ``formação e aprendizagem de toda a vida cristã''. Os catecúmenos ``já estão unidos à Igreja'', já são ``da casa de Cristo'' e a Igreja os abraça com amor e solicitude. Onde o Batismo de crianças se tornou a forma habitual, mantém-se a necessidade de um catecumenato pós-batismal, que é o espaço próprio da catequese.

O sentido e a graça do Batismo aparecem claramente nos ritos da sua celebração. O sinal da cruz, no início, manifesta a marca de Cristo naquele que passa a pertencer-Lhe; o anúncio da Palavra ilumina os candidatos e suscita a resposta da fé, pois o Batismo é, de modo particular, o ``sacramento da fé''. Pronuncia-se sobre o candidato um ou vários exorcismos, expressão da libertação do pecado e do demônio, e ele renuncia expressamente a Satanás antes de professar a fé da Igreja. A água batismal é então consagrada por uma oração de \emph{epiclese} --- a invocação ao Pai para que envie o Espírito Santo e o sacramento se realize ---, pela qual a Igreja pede que, pelo poder do Espírito Santo, os que nela forem batizados ``nasçam da água e do Espírito''.

O rito essencial é o Batismo propriamente dito: a tríplice imersão na água batismal, ou a tríplice infusão de água sobre a cabeça, acompanhada pelas palavras ``N., eu te batizo em nome do Pai e do Filho e do Espírito Santo''. Segue-se a unção com o santo Crisma, sinal do dom do Espírito Santo ao novo batizado, que se torna ``cristão'', isto é, ``ungido'' pelo Espírito, incorporado em Cristo sacerdote, profeta e rei. A veste branca simboliza que o batizado ``se revestiu de Cristo'' e ressuscitou com Ele; a vela acesa no Círio Pascal significa que Cristo iluminou o neófito. A primeira Comunhão eucarística, para os adultos batizados, completa a iniciação, introduzindo o novo cristão no ``banquete das núpcias do Cordeiro''.

\section{Quem pode receber o Batismo (1246--1255)}

Todo o ser humano ainda não batizado pode e deve receber o Batismo. Desde os princípios da Igreja, o Batismo dos adultos é a situação mais corrente nas terras de missão, onde o anúncio do Evangelho ainda é recente. O catecumenato tem nesse contexto um lugar central: dispõe os candidatos para o acolhimento do dom de Deus no Batismo, Confirmação e Eucaristia, iniciando-os no mistério da salvação, na prática dos costumes evangélicos e na vida de fé, liturgia e caridade do povo de Deus. A conversão que o catecumenato exige não é uma conquista humana, mas uma resposta à iniciativa divina da graça.

O Batismo das crianças é também uma prática apostólica imemorial da Igreja. Nascidas com uma natureza humana decaída e manchada pelo pecado original, as crianças têm necessidade do novo nascimento no Batismo para serem libertas do poder das trevas e transferidas para o domínio da liberdade dos filhos de Deus. A pura gratuidade da graça da salvação é particularmente manifesta no Batismo das crianças, pois elas não podem apresentar mérito algum. Por isso, a Igreja e os pais privariam a criança da graça inestimável de se tornar filha de Deus se não lhe conferissem o Batismo pouco depois do nascimento.

O Batismo é o sacramento da fé, mas a fé tem necessidade da comunidade dos crentes. Só na fé da Igreja é que cada fiel pode crer: a fé que se requer para o Batismo não é uma fé perfeita e amadurecida, mas um princípio chamado a desenvolver-se. Por isso, ao catecúmeno ou ao seu padrinho se pergunta: ``Que pedis à Igreja de Deus?'' --- e a resposta é: ``A fé!''. Em todos os batizados, crianças ou adultos, a fé deve crescer depois do Batismo; é por isso que a Igreja celebra todos os anos, na Vigília Pascal, a renovação das promessas batismais.

Para que a graça batismal possa desenvolver-se, é fundamental a ajuda dos pais e dos padrinhos. O padrinho ou a madrinha devem ser pessoas de fé sólida, capazes e preparadas para ajudar o novo batizado no caminho da vida cristã; o seu múnus é um verdadeiro ofício eclesial. Toda a comunidade eclesial tem uma parte de responsabilidade no desenvolvimento e na guarda da graça recebida no Batismo. A preparação para o Batismo conduz apenas ao umbral da vida nova; o restante do percurso é acompanhado pela catequese, pela oração e pelo testemunho fraterno da Igreja.

\section{O ministro do Batismo (1256)}

São ministros ordinários do Batismo o bispo e o presbítero e, na Igreja latina, também o diácono. Em caso de necessidade --- perigo de morte, ausência de ministro ordenado ---, qualquer pessoa, mesmo não batizada, pode batizar, desde que tenha a intenção de fazer o que faz a Igreja e empregue a forma trinitária: ``Eu te batizo em nome do Pai e do Filho e do Espírito Santo''. A Igreja vê a razão desta possibilidade na vontade salvífica universal de Deus e na necessidade do Batismo para a salvação. A intenção requerida é a de querer fazer o que faz a Igreja quando batiza, mesmo que o ministro não compreenda plenamente o mistério que está realizando.

\section{A necessidade do Batismo (1257--1261)}

O próprio Senhor afirma que o Batismo é necessário para a salvação: ``Se alguém não nascer da água e do Espírito, não pode entrar no Reino de Deus''. Por isso, ordenou aos seus discípulos que anunciassem o Evangelho e batizassem todas as nações. O Batismo é necessário para a salvação de todos aqueles a quem o Evangelho foi anunciado e que tiveram a possibilidade de pedir este sacramento. A Igreja não conhece outro meio garantido para assegurar a entrada na bem-aventurança eterna; por isso, tem cuidado em não negligenciar a missão que recebeu do Senhor de fazer ``renascer da água e do Espírito'' todos os que podem ser batizados. Ao mesmo tempo, Deus, que ligou a salvação ao sacramento do Batismo, não está Ele próprio prisioneiro dos seus sacramentos.

Desde sempre, a Igreja reconhece que aqueles que morrem por causa da fé sem terem recebido o Batismo são batizados pela sua morte por Cristo e com Cristo. Este ``Batismo de sangue'' produz os frutos do Batismo, mesmo sem ser sacramento. Para os catecúmenos que morrem antes de ser batizados, o desejo explícito de o receber, unido ao arrependimento dos seus pecados e à caridade, garante-lhes a salvação que não puderam receber pelo sacramento. Este ``Batismo de desejo'' exprime a profunda primazia da graça divina em relação ao rito sacramental.

Além do batismo de sangue e do batismo de desejo, a Igreja afirma que todo o homem que, na ignorância do Evangelho de Cristo e da sua Igreja, procura sinceramente a verdade e faz a vontade de Deus conforme o conhecimento que dela tem, pode salvar-se, associado ao mistério pascal ``por um modo só de Deus conhecido''. Podemos supor que tais pessoas teriam desejado explicitamente o Batismo se dele tivessem conhecido a necessidade. Isto não relativiza a urgência da missão, mas revela a amplitude da misericórdia divina: Deus quer que todos os homens se salvem e dá a cada um os meios necessários para alcançar a salvação.

Quanto às crianças que morrem sem Batismo, a Igreja não pode senão confiá-las à misericórdia de Deus, como o faz no rito do respectivo funeral. A grande misericórdia de Deus, ``que quer que todos os homens se salvem'', e a ternura de Jesus para com as crianças, que O levou a dizer ``Deixai vir a Mim as criancinhas'', permitem-nos esperar que haja um caminho de salvação para elas. Por isso, é ainda mais premente o apelo da Igreja a que não se impeçam as criancinhas de virem a Cristo, pelo dom do santo Batismo. Esta esperança não elimina a urgência da prática batismal, mas manifesta que o amor de Deus não pode ser encerrado em fronteiras que nós próprios estabeleçamos.

\section{A graça do Batismo (1262--1274)}

Os efeitos do Batismo são significados pelos elementos sensíveis do rito: a imersão na água evoca os simbolismos da morte e da purificação, mas também da regeneração e da renovação. Os dois efeitos principais são a purificação dos pecados e o novo nascimento no Espírito Santo. Pelo Batismo, todos os pecados são perdoados --- o pecado original e todos os pecados pessoais ---, bem como todas as penas devidas ao pecado. Naqueles que foram regenerados, nada resta que os possa impedir de entrar no Reino de Deus: nem o pecado de Adão, nem o pecado pessoal, nem as suas consequências.

No batizado permanecem, porém, certas consequências temporais do pecado: os sofrimentos, a doença, a morte, as fragilidades de caráter e, sobretudo, uma inclinação para o pecado que a Tradição chama concupiscência ou ``isca do pecado''. Esta inclinação foi deixada ``para os nossos combates'' e ``não pode fazer mal àqueles que, não consentindo nela, resistem corajosamente pela graça de Cristo''. O Batismo não somente purifica de todos os pecados, mas também faz do neófito uma ``nova criatura'': filho adotivo de Deus, participante da natureza divina, membro de Cristo, coerdeiro com Ele e templo do Espírito Santo.

A Santíssima Trindade confere ao batizado a graça santificante, a graça da justificação, que o torna capaz de crer em Deus, esperar n'Ele e O amar pelas virtudes teologais; lhe dá o poder de viver e agir sob a moção do Espírito Santo e pelos seus dons; e lhe permite crescer no bem pelas virtudes morais. O Batismo incorpora na Igreja: das fontes batismais nasce o único Povo de Deus da Nova Aliança, que ultrapassa todos os limites naturais das nações, culturas, raças e sexos. Os batizados tornam-se ``pedras vivas'' para a edificação espiritual e participam no sacerdócio de Cristo, na sua missão profética e real, sendo ``raça eleita, sacerdócio de reis, nação santa, povo que Deus tornou seu''.

O Batismo marca o cristão com um selo espiritual indelével --- o ``caráter'' --- da sua pertença a Cristo. Esta marca não é apagada por nenhum pecado e, por isso, o Batismo não pode ser repetido. O selo batismal capacita e compromete os cristãos a servir a Deus mediante uma participação viva na santa liturgia da Igreja e a exercer o sacerdócio batismal pelo testemunho duma vida santa e duma caridade eficaz. É também o vínculo sacramental da unidade entre todos os cristãos: aqueles que creem em Cristo e foram devidamente batizados estão numa certa comunhão --- embora não perfeita --- com a Igreja Católica, e são justamente reconhecidos pelos filhos da Igreja como irmãos no Senhor.

\section{Em síntese (1275--1284)}

A iniciação cristã realiza-se pelo conjunto de três sacramentos: o Batismo, que é o princípio da vida nova; a Confirmação, que é a consolidação dessa vida; e a Eucaristia, que alimenta o discípulo com o corpo e o sangue de Cristo em vista da sua transformação n'Ele. Segundo a vontade do Senhor, o Batismo é necessário para a salvação, como o é a própria Igreja, na qual o Batismo introduz. O rito essencial consiste em mergulhar na água o candidato ou em derramar água sobre a sua cabeça, pronunciando a invocação da Santíssima Trindade.

O fruto do Batismo, ou graça batismal, é uma realidade rica que inclui: a remissão do pecado original e de todos os pecados pessoais; o renascimento para uma vida nova, pela qual o homem se torna filho adotivo do Pai, membro de Cristo e templo do Espírito Santo; e a incorporação na Igreja, corpo de Cristo, tornando-se o batizado participante do sacerdócio de Cristo. O Batismo imprime na alma um sinal espiritual indelével --- o caráter --- que consagra o batizado para o culto da religião cristã; por causa dessa marca indelével, o Batismo não pode ser repetido.

Os que sofrem a morte por causa da fé, os catecúmenos e todos aqueles que, sob o impulso da graça, sem conhecerem a Igreja, procuram sinceramente a Deus e se esforçam por cumprir a sua vontade, podem salvar-se mesmo sem terem recebido o Batismo. Desde os tempos mais antigos, o Batismo é administrado também às crianças, visto ser uma graça e um dom de Deus que não supõem méritos humanos; as crianças são batizadas na fé da Igreja, e a entrada na vida cristã lhes dá acesso à verdadeira liberdade. Quanto às crianças que morrem sem Batismo, a liturgia da Igreja convida-nos a ter confiança na misericórdia divina e a rezar pela sua salvação.

Em caso de necessidade, qualquer pessoa pode batizar, desde que tenha a intenção de fazer o que a Igreja faz e derrame água sobre a cabeça do candidato, dizendo: ``Eu te batizo em nome do Pai e do Filho e do Espírito Santo''. Este dado revela a amplitude da solicitude da Igreja pela salvação de todos: o Batismo não é propriedade de nenhuma casta sacerdotal, mas dom oferecido a todos por meio de qualquer membro da família humana que queira cooperar com a graça divina. Recebendo o Batismo, o cristão não inaugura apenas uma vida nova individual: entra na comunhão dos batizados de todos os tempos e lugares, incorporado no mistério do Corpo de Cristo que percorre a história em direção ao Reino definitivo.
